% Generated by tools/supp_refs.py from main.aux; do not edit.
\makeatletter
\newlabel{sec:intro}{{1}{1}{Introduction}{Doc-Start}{}}
\newlabel{sec:regions}{{2}{4}{Trimmed regions and the witness}{Doc-Start}{}}
\newlabel{eq:trimset}{{2.1}{4}{Trimmings and total variation}{Doc-Start}{}}
\newlabel{eq:tvintersect}{{2.2}{4}{Trimmings and total variation}{Doc-Start}{}}
\newlabel{eq:contam}{{2.3}{4}{Trimmings and total variation}{Doc-Start}{}}
\newlabel{ass:regularity}{{2.1}{5}{}{Doc-Start}{}}
\newlabel{eq:trimMMDpop}{{2.4}{5}{Trimmed regions}{Doc-Start}{}}
\newlabel{eq:region}{{2.5}{5}{Trimmed regions}{Doc-Start}{}}
\newlabel{prop:georep}{{2.2}{5}{Geometric representation}{Doc-Start}{}}
\newlabel{eq:setdist}{{2.6}{5}{Geometric representation}{Doc-Start}{}}
\newlabel{lem:attainment}{{2.3}{5}{Compactness and attainment}{Doc-Start}{}}
\newlabel{fig:regions}{{1}{6}{Left: the regions $\mathcal C_\alpha (P)$ and $\mathcal C_\alpha (Q)$ of laws with weights $(\regWeightsP )$ and $(\regWeightsQ )$ on three points whose feature vectors are unit vectors at $120$ degrees in the plane, drawn isometrically (untrimmed embeddings as dots), at $\alpha =\regAlphaLo $, at $\alpha =\TV (P,Q)=\regTV $ and at $\alpha =\regAlphaHi $. Right: the distance $D_\alpha (P,Q)$, computed exactly, with these levels marked}{Doc-Start}{}}
\newlabel{eq:zeroeq}{{2.7}{6}{Trimmed regions}{Doc-Start}{}}
\newlabel{lem:cvar}{{2.4}{6}{Support function}{Doc-Start}{}}
\newlabel{eq:cvar}{{2.8}{6}{Support function}{Doc-Start}{}}
\newlabel{eq:hausdorff}{{2.9}{7}{Trimmed regions}{Doc-Start}{}}
\newlabel{sec:statistic}{{2.3}{7}{The statistic, its witness and the trimming it selects}{Doc-Start}{}}
\newlabel{eq:qp}{{2.10}{7}{The statistic, its witness and the trimming it selects}{Doc-Start}{}}
\newlabel{eq:witness}{{2.11}{7}{The statistic, its witness and the trimming it selects}{Doc-Start}{}}
\newlabel{prop:bangbang}{{2.5}{7}{Structure of the empirical trimming}{Doc-Start}{}}
\newlabel{sec:upper}{{2.4}{8}{No informative upper bound}{Doc-Start}{}}
\newlabel{thm:impossible}{{2.6}{8}{No informative upper bound, even for bounded densities}{Doc-Start}{}}
\newlabel{sec:analytic}{{3}{8}{Testing a trimmed null, and an analytic baseline}{Doc-Start}{}}
\newlabel{eq:null}{{3.1}{8}{Testing a trimmed null, and an analytic baseline}{Doc-Start}{}}
\newlabel{thm:plugin}{{3.1}{9}{Plug-in threshold}{Doc-Start}{}}
\newlabel{lem:oneevent}{{3.2}{9}{One event for all levels}{Doc-Start}{}}
\newlabel{eq:test}{{3.2}{9}{A threshold for the whole null}{Doc-Start}{}}
\newlabel{sec:alpha}{{3.2}{10}{The estimator}{Doc-Start}{}}
\newlabel{eq:alphaest}{{3.3}{10}{The estimator}{Doc-Start}{}}
\newlabel{sec:samplesize}{{3.3}{10}{How many observations}{Doc-Start}{}}
\newlabel{eq:samplesize}{{3.4}{10}{How many observations}{Doc-Start}{}}
\newlabel{sec:split}{{4}{10}{The split-sample lower bound}{Doc-Start}{}}
\newlabel{sec:ramp}{{4.1}{10}{A ramp of the witness}{Doc-Start}{}}
\newlabel{lem:ramp}{{4.1}{10}{Cut points turn the separation into a ramp}{Doc-Start}{}}
\newlabel{sec:betting}{{4.2}{11}{A betting test}{Doc-Start}{}}
\newlabel{thm:betting}{{4.2}{12}{Validity of the betting test}{Doc-Start}{}}
\newlabel{sec:lower}{{4.3}{12}{The lower confidence bound}{Doc-Start}{}}
\newlabel{cor:lower}{{4.3}{12}{Lower confidence bound for the contamination level}{Doc-Start}{}}
\newlabel{eq:downclosed}{{4.1}{12}{Lower confidence bound for the contamination level}{Doc-Start}{}}
\newlabel{prop:loconsistency}{{4.4}{13}{Consistency of the lower bound}{Doc-Start}{}}
\newlabel{sec:rates}{{5}{13}{Rates and optimality}{Doc-Start}{}}
\newlabel{sec:rate}{{5.1}{13}{A margin condition}{Doc-Start}{}}
\newlabel{ass:margin}{{5.1}{13}{Polynomial margin}{Doc-Start}{}}
\newlabel{ex:separated}{{5.2}{14}{Separated contamination}{Doc-Start}{}}
\newlabel{sec:linearmargin}{{5.2}{14}{Linear and slower margins}{Doc-Start}{}}
\newlabel{prop:newslope}{{5.3}{14}{Touching contaminations}{Doc-Start}{}}
\newlabel{prop:newapprox}{{5.4}{14}{Margin from an approximate witness}{Doc-Start}{}}
\newlabel{prop:newsobolev}{{5.5}{15}{Sobolev kernels}{Doc-Start}{}}
\newlabel{prop:newgauss}{{5.6}{15}{Gaussian kernel}{Doc-Start}{}}
\newlabel{sec:bothrates}{{5.3}{15}{The rate, and when it is optimal}{Doc-Start}{}}
\newlabel{sec:optimality}{{5.3}{15}{The rate, and when it is optimal}{Doc-Start}{}}
\newlabel{thm:loratemain}{{5.7}{15}{Rate of the split-sample lower bound}{Doc-Start}{}}
\newlabel{prop:lower2pt}{{5.8}{16}{Two-point lower bound}{Doc-Start}{}}
\newlabel{cor:minimax}{{5.9}{16}{Minimax rate}{Doc-Start}{}}
\newlabel{fig:excess}{{2}{17}{The rescaled path $(1-\alpha )D_\alpha $ for $P=N(1,1)$ and $Q=N(-1,1)$, whose total variation is $\alpha ^\star =2\Phi (1)-1\approx \fxAlphaStar $, with the Laplace kernel (a) and the Gaussian kernel (b), both with $\sigma =1$ (dots: computed levels; line: monotone interpolation). At $\alpha =\fxLevelLap $ in (a) and $\alpha =\fxLevelGauss $ in (b) the dashed line of slope $-(t^\star -u^\star )$ supports the path; the excess of the population ramp over $\alpha $ is at least the distance (arrow) from $\alpha $ to the zero of the line (tick), by Proposition~\ref {prop:excess}, and the inset of (b) enlarges the region near $\alpha ^\star $. That distance is a fraction $\fxFracLap $ of $\alpha ^\star -\alpha $ for the Laplace kernel and $\fxFracGauss $ for the Gaussian kernel at this level, closer to $\alpha ^\star $. The Gaussian fraction, $\fxFracGaussAtLap $ at $\alpha =\fxLevelGaussAtLap $, falls toward zero as $\alpha \uparrow \alpha ^\star $, where its path is flatter than every power}{Doc-Start}{}}
\newlabel{sec:touching}{{5.4}{17}{Touching contaminations, and why the margin is pessimistic}{Doc-Start}{}}
\newlabel{prop:excess}{{5.10}{18}{Excess of the population ramp}{Doc-Start}{}}
\newlabel{thm:laplacerate}{{5.11}{18}{Near-parametric rate with the Laplace kernel}{Doc-Start}{}}
\newlabel{eq:condR}{{5.1}{18}{Near-parametric rate with the Laplace kernel}{Doc-Start}{}}
\newlabel{sec:trimming}{{6}{19}{The trimming}{Doc-Start}{}}
\newlabel{eq:popdensity}{{6.1}{19}{A parametric rate at a fixed level}{Doc-Start}{}}
\newlabel{eq:qdens}{{{Q}}{19}{A parametric rate at a fixed level}{Doc-Start}{}}
\newlabel{eq:qdensY}{{{Q$'$}}{19}{A parametric rate at a fixed level}{Doc-Start}{}}
\newlabel{cor:trimrate}{{6.1}{20}{A parametric rate for the witness and the kept set}{Doc-Start}{}}
\newlabel{sec:trimest}{{6.2}{20}{The trimming at the estimated level}{Doc-Start}{}}
\newlabel{prop:alphahat}{{6.2}{20}{The trimmed sample at the data-driven level}{Doc-Start}{}}
\newlabel{cor:alphalo}{{6.3}{20}{The trimmed sample at the lower bound}{Doc-Start}{}}
\newlabel{sec:numerical}{{7}{21}{Numerical results}{Doc-Start}{}}
\newlabel{sec:numharder}{{7.1}{21}{Simulations}{Doc-Start}{}}
\newlabel{tab:hardern80}{{1}{22}{Sample size $n=m$ (mixed types: $n$, with $m=5n$) at which each method reaches power $0.8$ against $H_{0,0}$ at $\delta =0.05$, from $\mnRepsPerSize $ replicates per size, interpolated between the two sizes that bracket $0.8$ and rounded to two significant digits; ${\le }\mnNmin $, ${>}\mnNcap $: reached at the smallest size, not reached at the largest. Bold: the smallest among the four tests that certify a level of $\TV $; italic: the smaller of the last two, which test $P=Q$ only. $\widehat \alpha $: the betting test of Theorem~\ref {thm:betting} on the Laplace ramp; $\widehat \alpha _{\mathrm {an}}$: the analytic test \eqref {eq:test} with the Gaussian kernel; GBM, logistic: the betting test on the labels of gradient boosting and logistic regression trained once on half of the data, a handicap removed in Figure~\ref {fig:recovery}; perm.\ MMD: permutation test with the Gaussian kernel at the pilot bandwidth of the analytic test and $\abPermB $ relabelings; MMDAgg: the test of \citet {schrab2023aggregated}; the last two were run up to $n=800$ ($\fourdigits {\harderNtwo }$ on the narrow bump), so ${>}800$ means not reached there. The analytic and permutation tests take their bandwidth from a separate pilot draw, not counted in $n$ (Section~\ref {sec:numharder})}{Doc-Start}{}}
\newlabel{fig:recovery}{{3}{23}{How much of $\alpha ^\star $ each lower bound recovers at $n=m=800$ (mixed types: $n=400$, $m=5n$): the mean over $\mnRecReps $ replicates of the largest level $j\alpha ^\star /8$ it certifies divided by $\alpha ^\star $, with $0$ when it certifies none, on a grid that uses the known $\alpha ^\star $; $\alpha ^\star =0.2$ unless shown, $\delta =0.05$; classifiers: labels bet on the batches of the ramp}{Doc-Start}{}}
\newlabel{sec:data}{{7.2}{24}{Real data}{Doc-Start}{}}
\newlabel{tab:realdata}{{2}{25}{The two bounds on the four contrasts (the last three from NHANES), $\delta =0.05$, grid mesh $0.05$; $n,m$: group sizes, of which the analytic baseline uses $70\%$ after a $30\%$ pilot. Analytic baseline: $\widehat D_0$ and $\tau (0)$, its statistic and threshold at $\alpha =0$; $n_{\mathrm {fire}}$, the balanced size of \eqref {eq:samplesize}; $\alpha _{\mathrm {last}}$, the largest rejected grid level, and bisection, the largest level rejected when bisecting to width $0.0125$ (Section~\ref {sec:alpha}), both at most the infimum \eqref {eq:alphaest}; `--': none. Split-sample bound, Laplace kernel: the smallest bound over eight random splits and the level certified by the average of their final wealths, each $\delta $-level for this kernel. The three NHANES statements hold jointly with probability at least $1-3\delta $}{Doc-Start}{}}
\newlabel{sec:appcohort}{{7.3}{25}{From the contamination level to a cohort}{Doc-Start}{}}
\newlabel{sec:discussion}{{8}{26}{Final comments}{Doc-Start}{}}
\newlabel{mainend}{{8}{27}{Final comments}{Doc-Start}{}}
\makeatother
